\documentclass[11pt]{article}

\usepackage[margin=1in]{geometry}
\usepackage[utf8]{inputenc}
\usepackage[T1]{fontenc}
\usepackage{lmodern}
\usepackage{amsmath,amssymb,amsthm}
\usepackage{booktabs}
\usepackage{microtype}
\usepackage[numbers,sort&compress]{natbib}
\usepackage{xcolor}
\usepackage[colorlinks=true,linkcolor=blue!50!black,citecolor=blue!50!black,urlcolor=blue!50!black]{hyperref}

\newcommand{\draftnote}[1]{\textcolor{red!70!black}{[\textbf{Draft:} #1]}}

\title{Learning Certified Modular Division from Answers:\\Training on a Cover of the Target
Machine, and Collapse Past a Solution\\[0.5em]
\large First draft, 15 October 2026: not for circulation}

\author{
  David Barkhuizen\\
  Independent researcher
  \and
  Claude (Opus 5.5)\\
  Anthropic
}

\date{}

\begin{document}

\maketitle

\draftnote{Companion to \texttt{snapped-window-recurrence.tex} (the design, the certificate, the
exclusion bound; cited as \citep{barkhuizen-2026-snapped-window-recurrence}), split from it on 15
October 2026. Authorship and publication route: as there. Waiting on: the registered comparison
(Section~\ref{sec:claim}, running); the literature search is done, its papers in the literature
repository.}

\begin{abstract}
Modular division, the running quotient $A \cdot B^{-1} \bmod p$ of two numbers read digit by digit,
is a finite transducer whose state is the pair of residues: 25 states at $p = 5$, 49 at $p = 7$. A
window-recurrent transformer with a snapped state, which learns and certifiably computes the
3-state carry automaton and modular sums of up to 5 states, failed to learn division in all 160
runs trained on the quotient alone. We report what it took to learn it, every network checked by a
complete certificate that proves it correct at every length. Three things changed. First, the
quotient alone gives each residue only a sparse, binary signal (whether it is zero); one network
also answering the running sum and difference, each asked for by a query token on the same input,
gets dense signals: the two counters are homomorphic images of the target machine, and together
they cover it. Second, a state factored into one codebook per
residue, with each codebook's next code a function of its own code and the column only, learned
nothing, and the same constraint blocks even a single modular counter that is otherwise learned in
minutes; letting each codebook read the window's computation lifts the block. Third, every run
that reached or nearly reached a solution and was trained on collapsed, so each curriculum step is
certified and training stops at the first proof. With these, 11 networks certify against the
25-state machine in screens of 2--8 seeds (5 of 8 with two codebooks, 4 of 4 with one, against 0 of
6 for the quotient alone), and one of 4 against the 49-state machine; each passes every long-test case to 1{,}024 columns
(375{,}000 a network at $p = 5$, 447{,}000 at $p = 7$). In a registered comparison at 20 fresh seeds an arm, the cover certifies 13 of 20 against the quotient alone's 1 of 20 (one-sided Fisher's exact $p = 6.9 \times 10^{-5}$), and all 14 certified networks pass every long-test case. The learned states
are not factored by residue: each certified network spreads the pair over both codebooks and
33--88 codes, which minimise to exactly the $p^2$ states.
\end{abstract}

\section{Introduction}

A companion paper \citep{barkhuizen-2026-snapped-window-recurrence} studies a small transformer
that reads a window of digit columns at a time and hands a state to the next window through a
gated recurrent cell whose output is snapped to the nearest of $m$ learned codes. Because the state
is snapped, a trained network is exactly a finite-state transducer, and a complete certificate
proves it correct at every length or yields a counterexample. That paper's networks learn the
3-state carry automaton of base-3 addition and modular sums of 4 and 5 states.

Modular division is the next rung and a different kind of machine. Its state is a product of two
independent counters, the residues of the two numbers read so far, and its output, the quotient,
depends on both together. The same design, trained on the quotient with codebooks of up to 10 times
the 25 states, failed in all 160 runs (Section~\ref{sec:failures}). This paper reports how the
failure was taken apart and what certified division at $p = 5$ and once at $p = 7$.

\paragraph{Contributions.} Each with the nearest prior work we found; ``not found'' is the result of
our search, not a claim that no such work exists.
\begin{enumerate}
\item \textbf{Certified product machines learned from answers.} Networks that certifiably compute
  the 25-state and 49-state division transducers, trained only on answers at every column, never
  on states or intermediate steps (contrast the expert traces of
  \citet{papazov-2025-exact-learning-arithmetic-differentiable-agents}); the minimal machines
  extracted by \citet{wang-2019-state-regularized-rnns} and
  \citet{oliva-2020-dual-recurrent-memory-processing} have at most a few states.
\item \textbf{Training on a cover of the target machine} (Section~\ref{sec:family}): beside the
  target's output, the outputs of its homomorphic images (here two modular counters, which together
  cover it), answered by query tokens on the same input, so that each coordinate of the state has a
  dense signal. Screened at 5 of 8 against 0 of 6 for the quotient alone, and registered at 20 fresh
  seeds an arm: 13 of 20 against 1 of 20 ($p = 6.9 \times 10^{-5}$). Nearest prior work: multitask learning's
  ``eavesdropping'' on a feature another task makes easy \citep{caruana-1997-multitask-learning};
  auxiliary supervision of a DFA's states \citep{dan-2022-robust-generalization-regular-languages},
  which with the full cover ours approaches, since the cover's outputs determine the state; we found
  no training on homomorphic images of a target machine, nor a case where a single approximation
  without a cover suffices.
\item \textbf{An enforced factorisation that blocks learning} (Section~\ref{sec:d10}): a codebook
  per counter whose next code may depend only on its own code and the column cannot learn even one
  counter; the factors that learn are coupled, and the machine they learn is not factored.
\item \textbf{Collapse past a solution, at scale} (Section~\ref{sec:collapse}): every run that
  reached or neared a solution and trained on lost it, including a run right on every sampled answer
  for four steps; certifying each step and stopping at the first proof produced every certified
  network.
\end{enumerate}

\section{Related work}
\label{sec:related}

The design, the certificate, discretised recurrent states and automaton extraction are reviewed in
the companion paper \citep{barkhuizen-2026-snapped-window-recurrence}. The algebraic terms of
Section~\ref{sec:family} (homomorphic image, or approximation; cover, or emulation; coordinates) are
those of the Krohn--Rhodes and holonomy decompositions as
\citet{egri-nagy-2015-computational-holonomy-decomposition,egri-nagy-2024-covering-lemma-transformation-semigroups}
use them; $k$-distinguishability is \citet{chow-1978-testing-software-design-finite-state-machines}'s.

\paragraph{Learned modular and recurrent arithmetic.} Learned arithmetic in recurrent models has been
addition and its relatives
\citep{kaiser-2015-neural-gpu,price-2016-neural-gpu-extensions-limitations,mohri-2026-rational-transductors};
\citet{papazov-2025-exact-learning-arithmetic-differentiable-agents} learn exact multiplication with
finite-state transducers supervised on an expert's steps. In
\citet{deletang-2022-neural-networks-chomsky-hierarchy} a 5-state modular expression task is regular,
and recurrent networks and LSTMs generalise on it where their transformer does not.
\citet{power-2022-grokking-small-algorithmic-datasets} include $x / y \bmod p$ among their grokking
tasks, a table of single symbols rather than a transducer over digits;
\citet{nanda-2023-progress-measures-grokking} reverse-engineer the learned algorithm for modular
addition. \citet{he-2024-learning-to-grok-modular-arithmetic} train one transformer on many linear
modular functions $a x + b y \bmod p$ (our sum and difference are the cases $(1, 1)$ and $(1, -1)$,
applied to running residues) and find that out-of-distribution generalisation in deeper models is
transient and needs early stopping. \citet{robinson-2020-extrapolating-algorithmic-learning} count a
network's computational states; \citet{buitrago-2025-length-generalization-recurrent-models}
attribute recurrent failures to states training never visits.
\citet{allchin-2026-reward-teach-state-hidden-automaton} find that reinforcement-learning agents
rarely learn the state of a permutation automaton, one in which no input maps two states to one
successor; division's target machine is one (each column permutes the $p^2$ pairs), as is the
modular sum, which the companion paper's networks learn.

\paragraph{Auxiliary outputs and intermediate supervision.} Multitask learning uses related tasks'
training signals as an inductive bias; among its mechanisms, a net can ``eavesdrop'' on a feature
that another task makes easy to learn and its own task makes hard
\citep{caruana-1997-multitask-learning}, which is our reading of the cover: a residue is hard to
learn from the quotient, whose signal about it is only its zero test, and easy from a counter.
Auxiliary tasks shape shared representations in reinforcement learning
\citep{jaderberg-2017-unsupervised-auxiliary-tasks}. Closest to ours,
\citet{dan-2022-robust-generalization-regular-languages} supervise a recurrent network's states with
the target DFA's (``state sequence auxiliary supervision'') and find it generalises far better on
modulo languages; with the full cover our outputs determine the state up to relabelling, so ours is
close to theirs, while the quotient with one approximation is not. Their label-free alternative, a
capped count of zeros, records more than the parity it serves; our approximations record less than
the state. Algorithmic reasoning with graph networks supervises every intermediate step of an
algorithm (``hints'') and finds positive transfer between related algorithms
\citep{velickovic-2020-neural-execution-graph-algorithms,velickovic-2022-clrs-algorithmic-reasoning-benchmark,ibarz-2022-generalist-neural-algorithmic-learner};
\citet{wies-2023-sub-task-decomposition-sequence-tasks} prove that intermediate supervision of
sub-task results can make an unlearnable composite learnable. Ours supervises no steps: every output
is a function of the input read so far, defined at every column whatever the network computes.

\paragraph{Discrete codes: factoring and collapse.} A codebook that is the Cartesian product of
smaller ones is product quantisation \citep{jegou-2011-product-quantization-nearest-neighbor}; our
factored state is one, each factor with its own recurrent cell. Decomposing a vector quantiser into
several smaller ones was introduced partly against \emph{index collapse}, in which only a few codes
are ever trained \citep{kaiser-2018-fast-decoding-discrete-latent-variables}; collapse is commonly
countered by moving averages, re-initialising or restarting unused codes
\citep{lancucki-2020-robust-training-vq-bottleneck,dhariwal-2020-jukebox,zheng-2023-online-clustered-codebook},
and \citet{huh-2023-straightening-straight-through-estimator} trace it to a drift of the embeddings
away from the codebook, with codes that lose their assignments no longer trained, and to the
asymmetry of the commitment loss. Separately, adaptive optimisers show cyclic instabilities long
after a network fits its training data, the norm of the last layer growing until training examples
are misclassified \citep{thilak-2022-slingshot-mechanism-adaptive-optimizers}.

\section{Tasks, network and certificate}
\label{sec:setup}

\paragraph{Division.} Two base-3 numbers $A$ and $B$ are read most significant digit first, a column's
token $3a_i + b_i$ (9 values). The output at every column is the running $(A \cdot B^{-1}) \bmod p$ of
the prefixes read, or ``undefined'' while $B \equiv 0$: $p + 1$ classes. Its state is the pair
$(A \bmod p, B \bmod p)$, each residue updated by $r' = (3r + d) \bmod p$, and its minimal automaton
has exactly $p^2$ states (25 and 49 for $p = 5, 7$), checked by partition refinement.

\paragraph{The target machine and its cover.}
\label{sec:family}
We call division's 25-state transducer the \emph{target machine}; the quotient is its output
function. We also train on two further outputs over the same input: the sum $(A + B) \bmod p$ and the
difference $(A - B) \bmod p$. A window of one column becomes three tokens, $[\,x, Q_{+}, Q_{-}\,]$:
the column's token answers the quotient, as before, and each query token's output is its readout,
through the one output layer. In the terms of
\citet{egri-nagy-2024-covering-lemma-transformation-semigroups}: the partition of the target's states
by $A + B$ is a congruence (every transition respects it), so the sum is a homomorphic image of the
target machine, a surjective morphism they call an \emph{approximation}; it is a modular counter of
its own, the companion paper's modular sum. So is the difference. For odd $p$ the two determine the
pair, $A = (S + D)/2$ and $B = (S - D)/2$, so the target machine embeds in the product of the two
counters: the product \emph{covers} (emulates) it, an injective morphism, and $S$ and $D$ are its
\emph{coordinates}. The quotient is not an approximation (its partition is not a congruence), and
neither is the product $A B \bmod p$.

Every set of outputs we train on identifies the state equally fast: with the quotient alone, the
quotient and the sum, or the sum and the difference as outputs, the target machine is
1-distinguishable \citep{chow-1978-testing-software-design-finite-state-machines}, every pair of
states told apart by some single column. What sets the cover apart is that its outputs determine
the current state outright (the output function is injective on states), so a reader may count it as
supervision of the state, and we do too. The quotient with the sum is neither: one approximation, not
a cover, its outputs not injective, and it also learns (Section~\ref{sec:isolate}). Each of these
facts was checked exhaustively at $p = 5$ and 7. The quotient alone gives a residue only a sparse,
binary signal: its information about each residue is its information about whether that residue is
zero. Only outputs are labelled, never states or intermediate steps; choosing the outputs is a
curriculum choice, made as the ladder of modular sums before division was. The code and its records
call the cover's arm \texttt{family}.

\paragraph{The network.} As in the companion paper, at $k = 1$ column a window: the window's tokens
embedded with a learned position, a causal transformer of 3 blocks (width 16, 2 heads, feed-forward
32) with a projection of the state added before every block, an output softmax per token, and a GRU
whose output is snapped to the nearest of $m$ learned codes, straight-through, with VQ-VAE's codebook
and commitment losses \citep{van-den-oord-2017-vq-vae-neural-discrete-representation}. New here, the
state may be \emph{factored} into $G$ factors, each with its own learned pooling query over the
window's tokens, its own GRU cell and its own codebook; the factors' codes together are the state,
up to $m_1 \cdots m_G$ joint states, and every block reads all of them. A pooling reads either the
tokens of the last block, which have seen the whole state (\emph{reading the output}, the companion
paper's network), or the embedded tokens before any state is added (\emph{reading before the
state}): then a factor's next code depends on its own code and the column alone, the independence
division's two counters have. The networks have 12{,}310 (two factors of 16 codes, $p = 5$),
12{,}344 ($p = 7$) and 17{,}414 parameters (one codebook of 128 codes, width 32), codebooks
included.

\paragraph{The certificate} is the companion paper's: every window from every reachable joint
code, mapped onto the task's automaton, every output and next state checked, the learned machine
minimised, the snap ratio and logit lead of every transition recorded. For the cover's arm the machine's
input is the column token, its queries attached, and the quotient's slot is read. \emph{Certified}
uses the network's own snaps; \emph{robustly certified} branches every near tie (snap ratio below
1.05) both ways. Every network called certified here passes both. The long tests run the uniform
and stratified catalogues at 128, 256 and 1{,}024 columns: 375{,}000 cases a network.

\paragraph{Feasibility and the exclusion bound.} Every cell can represent its target (two factors
of 16 codes give 256 joint states for 25 or 49; one codebook of 128 for 25), and the curriculum's
cap is the exclusion bound $m + q - 1$ of the companion paper with $m$ the joint state count: 280,
152 and 304 columns. Before the $p = 7$ probe the cap was a fixed 280, below $p = 7$'s 304 and above
the one-codebook cells' 152; no run reached step 34, so no result depended on it. No run approaches
the bound: a certificate, not the training length, excludes memorising.

\paragraph{Training.} The curriculum: step $j$ trains on lengths $1, \dots, j$ columns, warm-started,
and ends when its check passes, after 10 epochs without improvement, or at 60 epochs; epochs of
$2^{15}$ examples, batches of 64, Adam at 0.003. At each step's end the quotient's machine is
certified, and a run stops at its first certified step, its network saved. Screens run 60 minutes a
run; the registered comparison runs 400 epochs. Runs are single-threaded on two CPUs (Intel
i7-9700K, AMD Ryzen 7 3700U).

\section{Results}

\subsection{The quotient alone: 0 of 160}
\label{sec:failures}

Under the companion paper's harness, trained on the quotient: at $p = 5$ and 7, a
curriculum and an all-at-once arm, $k \in \{1, 4\}$, $m \in \{32, 64\}$, width 16, 5 seeds,
\textbf{0 of 80} work or certify. (The $p = 7$, $m = 32$ cells could not succeed, 32 codes against 49
states; they ran as registered.) At $k = 4$ no run fits its training lengths; at $k = 1$ every run is
a near miss at length 3, the first length that uses every transition. With 5 and 10 times as many
codes as states ($m$ = 128--512, $k = 1$, 10 seeds), again \textbf{0 of 80}: the certificate's map
covers all 25 states at $p = 5$ in every all-at-once $m = 256$ network, yet 79 of the 80
counterexamples are a wrong next state. More codes reduce the errors at length 3 (the best network,
4 of 729 pairs wrong) but none reaches zero. These runs predate the exclusion bound and the
stop-at-proof protocol: they trained far short of the bound, to completion.

\subsection{An enforced factorisation learns nothing}
\label{sec:d10}

Since division's state is two independent counters, we gave the network two codebooks reading
before the state, so that each could hold one counter and nothing else; existence tests, each
step certified and stopped at the first proof, the curriculum at $k = 1$. Two factors of 6 codes
(C9, 6 seeds, stopped at 4.7 hours at steps 19--21), two of 16 (C9$'$, 6 seeds, stopped at steps
9--12) and three of 16 trained in turns (C9$''$, 9 seeds, 4 hours): \textbf{none certified}. Their
diagnoses agree on a mechanism: trained together, the factors settle on the zero tests (each factor's
share of a residue at most 0.41, the nonzero quotients 31--44\% right against a guess of
25\%), and a factor trained alone in its turn collapses to one code.

A check on a task the design had learned isolates the constraint. On the modular sum mod 5, whose
single counter the companion paper's network certified in 6 of 6 runs, one factor reading before the
state passes steps 1--2 and then plateaus at every step: 0 of 2 certified in 30 minutes; the same
factor reading the output certifies as before, at steps 4 and 5: 2 of 2. Reading before the state, a
factor's cell sees only the embedded column and must compute the whole transition itself; reading
the output, the three blocks do most of it. (The explanation is a reading, not tested apart, and the
check has 2 seeds an arm.) Every factored division run had read before the state.

\subsection{\texorpdfstring{Training on the cover certifies division at $p = 5$}{Training on the cover certifies division at p = 5}}

\paragraph{The first probe} (registered with the poolings reading before the state, as C9): the
cover and the quotient alone, 2 seeds each, an hour: none works, the factors' shares of any
residue at most 0.18. \textbf{With the poolings reading the output} (a variation after the D10
check, not registered beforehand), \textbf{the cover certified 2 of 2 and the quotient alone 0
of 2.} The two networks were saved before the probe stopped at proofs (at steps 11 and 12); each
certifies against the 25-state machine in both verdicts, minimised to 25 from 32--53 codes reached,
logit leads at least 2.03, and their sums certify against the modular sum too. Each was right on
every answer at 3--12 columns and at 24 and 48 from steps 6--7. One run's step-6 network, 98.7\%
right at 48 columns on samples, fails the certificate with 1{,}025 wrong outputs.

\paragraph{Consolidation} (registered, 6 new seeds, stopping at the first proof): \textbf{3 of 6
certified}, at steps 7, 8 and 10, within 25 minutes; with the first two, \textbf{5 of 8} (Wilson
95\% interval 0.31--0.86). Each reaches 57--88 joint codes and minimises to the 25 states, with
13--35 near ties branched and logit leads at least 0.60. Every one of the five passes all 375{,}000
long-test cases.

\paragraph{Not factored.} At a solved step each factor uses 14 codes and holds 0.37--0.71 of every
direction $xA + yB$; no factor's code is a function of one residue, and the pair of codes takes 63
values for the 25 states. The certificate's readout agrees. The network solved division with its
state spread over both codebooks.

\subsection{What the cover and the codebooks contribute}
\label{sec:isolate}

\begin{table}[t]
\centering
\caption{Division screens, each run 60 minutes, stopped at its first certified step; every certified
network certified in both verdicts and passing all its long-test cases. $p = 5$ unless noted; the
poolings read the output.}
\label{tab:screens}
\begin{tabular}{llrl}
\toprule
outputs besides the quotient & state & certified & steps, minutes \\
\midrule
sum, difference (cover) & two factors of 16 & 5 of 8 & 7--12, within 25 \\
none & two factors of 16 & 0 of 6 & -- \\
sum, difference (cover) & one codebook of 128 & 4 of 4 & 5--8, within 18 \\
sum (an approximation) & one codebook of 128 & 2 of 2 & 9--10, within 13 \\
cover, $p = 7$ & two factors of 16 & 1 of 4 & 7, 11 \\
\bottomrule
\end{tabular}
\end{table}

\begin{itemize}
\item \textbf{The quotient alone} (4 more seeds): 0 of 4, so 0 of 6 against the cover's 5 of 8
  (Fisher's exact, two-sided $p = 0.031$). The two arms ran as separate probes rather than
  interleaved seeds, so this is a screen; the registered comparison follows.
\item \textbf{One codebook} of 128 codes: 4 of 4 certified, within 18 minutes, against 5 of 8 for two
  factors ($p = 0.49$): the second codebook is not needed at $p = 5$. We keep two factors as the
  working configuration, since one codebook's size must grow with the product of the residues.
\item \textbf{The quotient and the sum}, without the difference, one codebook: 2 of 2, with the
  widest logit leads yet (at least 4.14). One approximation is enough, without a cover. Since the
  quotient alone is 1-distinguishable as well (Section~\ref{sec:family}), what the sum adds is not
  faster identification of the state but a dense signal on a counter that closes on itself.
\end{itemize}

\subsection{\texorpdfstring{$p = 7$: 49 states}{p = 7: 49 states}}

The cover with two factors of 16 (256 joint states, 5.2 times the 49), 4 seeds, an hour each:
\textbf{1 of 4 certified}, at step 7 after 11.2 minutes, against the 49-state machine in both
verdicts: 88 joint codes reached (87 without branching), minimised to 49, 34 near ties branched,
logit leads at least 2.05, snap ratios at least 1.004. It passes all 447{,}000 of its long-test cases at 128, 256 and 1{,}024
columns. The
quotient alone with one codebook had failed in all 80 runs at $p = 7$ (Section~\ref{sec:failures}).
The other three runs collapsed (minimised to 17, 5 and 20 states, 30--33\% right at 24 columns at
their peaks) and none was still rising when its hour ended.

\subsection{The registered comparison}
\label{sec:claim}

Registered before its data: the cover against the quotient alone at $p = 5$, two factors of 16,
reading the output, 20 fresh seeds an arm (100--119, none used in the screens), 400 epochs a run,
stop at certified, every certified network recertified in both verdicts and long-tested.
\textbf{Works} if the cover's count exceeds the quotient alone's, one-sided Fisher's exact
$p < 0.05$.

\textbf{It works.} The cover certified 13 of 20 (Wilson 95\% interval 0.43--0.82), the quotient
alone 1 of 20 (0.01--0.24); one-sided Fisher's exact $p = 6.9 \times 10^{-5}$. The cover's
successes certified at steps 6--12 (median 8) after 99--257 epochs (median 177), 10--55 minutes;
the quotient alone's one at step 10 after 202 epochs. Every certified network recertified in both
verdicts against the 25-state machine (35--74 joint codes reached, each minimised to 25), and every
one passes all 375{,}000 of its long-test cases, 5{,}250{,}000 in all. Every failure spent its 400
epochs (no run ended on the time guard).

The failures differ by arm. Four of the cover's seven came near a solution (every kind of answer at
least 95\% right on the sampled cases) and lost it; two of the quotient alone's nineteen did.
Large drops (the worst-answered kind from at least 0.9 right down by more than 0.3 within five
epochs) are common
and mostly recovered: 16 in 10 of the cover's runs, 6 of them in runs that certified, and 8 in 4 of
the quotient alone's; 17 of the 24 returned to at least 0.9. A failure is a drop not recovered from;
whether the codebooks collapsed did not decide it. The curriculum moves to the next length when a
step plateaus, and a drop leaves the step's best score at the peak before it, so a run that drops
has ten epochs to recover or goes on broken.

\subsection{Collapse past a solution}
\label{sec:collapse}

Every division run that reached or neared a solution and kept training lost it. In the first probe,
both cover runs collapsed after solving (at steps 12 and 14, one to 4 codes a factor), after their
networks had been saved; both quotient-alone runs collapsed too. In the consolidation, the three runs that
failed had each progressed and then collapsed, minimised to 6--11 states, 23--25\% right at 24
columns; one of them was right on every sampled answer at steps 6--9 (100\% at 24 columns, 99.99\%
at 48, minimised to 30 states) without certifying, and then collapsed to 2 states. At $p = 7$ all three
failures are collapses. Successes needed 73--215 epochs; no failure was slow progress, so a longer
budget would not raise the rate and stopping the collapse would. The companion paper saw the same on
modular sums \citep{barkhuizen-2026-snapped-window-recurrence}: certified networks destroyed by
further training.

A candidate mechanism is the commitment loss, which pulls each cell's output towards its code: a
cell whose state passes nothing on meets it most cheaply with a constant. In a probe on two
networks, a re-drawn factor with the commitment on fell from 11--12 codes to 2--8 within 5 epochs,
and with it off kept 14--16. That is suggestive, not shown, and it describes a factor starting, not
a solved network collapsing.

The collapse is index collapse in the sense of vector quantisation: the codes in use fall from tens
to a few a factor (1--5 where recorded), the state minimised to 2--20 states
\citep{kaiser-2018-fast-decoding-discrete-latent-variables}. It differs from the
case usually described in coming late, after a codebook had been fully in use and the task solved.
Two published accounts fit that timing: the drift of the cell's outputs away from the codes with
the commitment loss's asymmetry \citep{huh-2023-straightening-straight-through-estimator}, and the
late instabilities of Adam after a network fits its data
\citep{thilak-2022-slingshot-mechanism-adaptive-optimizers}; and transient generalising solutions
that need early stopping are reported for modular arithmetic in transformers
\citep{he-2024-learning-to-grok-modular-arithmetic}. We have not tested either account, nor the
remedies of the vector-quantisation literature (moving-average codebooks, code restarts, affine
re-parameterisation).

\section{Discussion}

\paragraph{What it took.} From 0 of 160 to certified networks, four things changed together: the
cover as further outputs, the poolings reading the output, the curriculum certified at each step and
stopped at the first proof, and the probe's shorter steps. The screens separate two of them: the
D10 check shows reading before the state blocks a counter the design otherwise learns, and the
cover against the quotient alone under the same settings is 5 of 8 against 0 of 6, the comparison
registered at 20 seeds. Stopping at the first proof is not separated on division; without it, every
run that solved would have ended collapsed.

\paragraph{Factored, but not by residue.} The design gave the network two codebooks for two
counters, and the network did not use them that way. A factorisation imposed by blocking the
factors' view of each other prevented learning; a free one learned a machine spread across both,
using the product codebook \citep{jegou-2011-product-quantization-nearest-neighbor} as one larger
code space.
Whether factored readouts can be had without blocking learning is open.

\paragraph{Limitations.}
\begin{itemize}
\item \textbf{Screens.} Every rate but the registered comparison's comes from 2--8 seeds; the
  reading-output variation was chosen after its registered probe failed.
\item \textbf{One rung up.} One network of 4 at $p = 7$; the rate there, and $p$ beyond 7, are open.
\item \textbf{The cover's outputs determine the state.} The quotient with the sum, whose outputs
  do not, also learned (2 of 2), so the state being readable is not needed; that the approximation
  is what helps is a hypothesis. Its test: the target with the product $A B \bmod p$, as fast to
  distinguish states and not an approximation.
\item \textbf{The collapse is described, not explained or prevented}; stopping at the first proof
  avoids it only where each step can be certified.
\end{itemize}

\section{Reproducibility}

Code: \texttt{davidbarkhuizen/indrajala-ml}. The probes, their settings, results, commits, machines
and SHA-256 manifests are recorded in \path{scripts/modular_family_probe.py}'s documentation and
\path{docs/division-workplan.md} (each probe registered there before its data), the factored-state
studies in \path{scripts/window_recurrent_study.py} and \path{docs/factored-state-workplan.md}.
Every certified network, with its certificates in both verdicts, long tests and training log, is
archived in \texttt{davidbarkhuizen/indrajala-models} (\path{window-recurrent/2026-10-15-division/}).
The $p = 7$ network and the registered comparison's 14 are there too.
The code (\texttt{davidbarkhuizen/indrajala-ml}, with its Rust crate
\texttt{davidbarkhuizen/indrajala-math-rust}) and the networks
(\texttt{davidbarkhuizen/indrajala-models}) are public, the paper's results at a tagged release;
the result files and their manifests are in the paper's datapack on Zenodo \draftnote{the release
tag and the datapack's DOI when deposited}.

\section*{Author contributions}
D.B. conceived the research programme and the design, made the design and experimental decisions,
set the experimental rules (feasibility of every cell, the exclusion bound, existence first,
probe-driven experiments, answers-only training) and directed the work. Claude (Opus 5.5,
Anthropic), working as an AI coding agent under D.B.'s direction, implemented the networks, probes
and certificate adapter, ran the experiments, carried out the analysis and drafted this paper.
\draftnote{D.B.\ to revise.}

\bibliographystyle{plainnat}
\bibliography{library,extra}

\end{document}
